\cotefile{156-4}
\scanpage{26}\mspage{25}
On reprend tout !

\uline{Définition-\unsure{scholie}} Une \uline{contrée} \inserted{$C$} est \unsure{definie|définie} par un ensemble $L$ (l'ens. des \uline{lieux} de la contrée), et la donnée de certains ensembles de parties de $L$, appelés les \uline{figures} \inserted{(de la contrée)}.

L'ens. $\unsure{\mathcal{F}|\mathscr{F}}$ des figures de $C$ est donc une \deleted{ensemble} partie de $P(P(L))$ (donc $F\in P(P(P(L)))$). Si \deleted{$\unsure{\mathcal{F}|F}$}$F$ est une figure, ses éléments\textsuperscript{\unsure{$\times$/}} (qui sont des parties de $L$) s'appellent les \uline{strates} (\uline{fermées}) de la figure $F$. \unsure{Là|Lá|La} réunion des strates \deleted{\unsure{fermées|ferm}} de $F$ est appelée le \uline{support} de $F$, \inserted{noté $|F|$.} On pose des \uline{axiomes des contrées}, qui vont être dégagés par la suite.

\uline{Cont}~1 \quad a) \unsure{La|la} partie vide de $P(L)$ est une \deleted{strate} figure, appelée ``\uline{figure vide}''. C'est donc la figure unique qui n'a pas de strates.

% slanted
\margin{\unsure{Contenu|Contenue} dans \uline{Cont} \unsure{4|5} plus bas}
\margin{$\big\Vert$}
b) $\forall\,x\in L$, \deleted{\unsure{la|le}}\inserted{il y a une} figure ayant \deleted{\unsure{la|les}} strate unique $\{x\}$. On l'appelle \uline{figure ponctuelle} \unsure{definie|définie} par $x$
\scanpage{27}\mspage{26}
\uline{Cont 2} \quad Soit $F$ une figure. \deleted{Alors} Pour tout $X\in F$, \deleted{\unsure{d'une|l'ens.}} \deleted{\unsure{soit}}\unsure{soient|soit}
\[
\partial X^{\deleted{\unsure{\circ|*}}}=\bigcup_{Y\in F_{<X}}Y,\qquad X^\circ=X\setminus\partial X
\]
(strate \uline{ouverte} de $F$ définie par $X$) Alors on a $\boxed{X^\circ\ne\varnothing}$. De plus, si $X,Y\in F$, alors
% slanted
\margin{Les \unsure{ens|ens.} de parties satisfaisant \unsure{ces|aux} conditions seront dits ``admissibles''}
\[
X\cap Y=\bigcup_{\substack{Z\in F\\\unsure{X|Z}\leq\unsure{x|X},Y}}F_{\unsure{Z|X}}
\]
i.e. \deleted{\unsure{tout|tous}} pour \unsure{tt|ttt} \deleted{\unsure{pt}} lieu\textsuperscript{\unsure{$x$}}/ qui appartient à deux strates fermées, \deleted{on a $\exists\,Z\in F$} il existe une \uline{strate} $Z\subset X\cap Y$ \deleted{telle que} \deleted{\unsure{le lieu}} qui contient $x$.

\uline{Corollaire} \deleted{\unsure{Si application|L'application} $F$} \inserted{\deleted{\unsure{par}} L'application,}
\[
\unsure{x|X}\longmapsto X^\circ\qquad F\longrightarrow P(|F|)
\]
\deleted{\unsure{i|in}}induit une bijection de $F$ avec l'ens. des classes d'une relation d'équivalence de $|F|$, i.e pour $X,Y\in F$, \deleted{tels que} \inserted{on a}
\[
\deleted{X\ne Y,\ \text{on a}\quad X^\circ\cap Y^\circ=\varnothing}
\]
\[
X\ne Y\iff X^\circ\cap Y^\circ=\varnothing
\]
(équivaut en fait à \uline{Cont 2})

\deleted{Définition} \inserted{\unsure{Cf.|Cf} \uline{Parties} \uline{saturées} de $|F|$ : p.ex. les $X\in F$}

\deleted{\uline{Définition} Soient $F,\unsure{G|C}$ deux figures. On dit qu'elles sont \uline{compatibles}, si \unsure{pour} les conditions équivalentes a) $|F|\cap|G|$ est saturé pour les relations d'\unsure{équivalence|équivalences} sur $|F|$, $|G|$, et la rel.}
\scanpage{28}\mspage{27}
\uline{Cont} 3 Si $F$ est une \deleted{\unsure{figaille|figille}} figure, et $F'$ une partie fermée de l'ens ordonné $F$ (ie. $X,Y\in F$, $Y\in F'$, $\deleted{X}\leq Y\implies X\in F$) alors $F'$ est une figure.

On dit que $F'$ est une sous-figure -- ainsi, les \unsure{sous figures|sous-figures} de $F$ sont en corr. biunivoque avec les parties fermées de l'ens ordonné $F$.

\unsure{Attention}, si $F',F$ sont deux figures, on peut avoir $F'\subset F$ sans que $F'$ soit une sous figure. (\unsure{Dans} le cas de l'exemple-type, prendre $F$ formé de $X_0\subset X_1$, $X_1$ une variété connexe, $F'=\{X_1\}$.)

% slanted
\margin{ici, $F'\subset F$ mais $F'\unsure{\nleq|\not\subseteq}F$ et $F',F$ ne sont \unsure{pas|$\ne$} compatibles, au sens plus bas. --}

On écrira \deleted{$\unsure{\leq|\subset}$} $F'\leq F$ pour $F'$ sous-figure de $F$. C'est une relation d'ordre entre figures.

\deleted{Soient $F,\unsure{G|C}$ deux figures. \deleted{\unsure{(plus gén. la défin.|(plus gén. la définition} \unsure{d'équivalence|d'équi}} \inserted{On dit que} $F$ et $G$ sont \deleted{\unsure{équivalentes}} \inserted{\unsure{ens\textsuperscript{t}.|ensemb.}} \inserted{compatibles}, si $F\cup G$ est}
\scanpage{29}\mspage{28}
Soient $F$ et $\unsure{G|C}$ \unsure{deux|des} figures. Conditions équivalentes (car $F$ et $G$ sont des ens. \uline{admissibles} de parties de $L$)

a) $|F|\cap|G|$ est une partie saturée de $|F|$ et de $|G|$, et les deux \unsure{rel,|rel.} d'équiv. induites sont égales

% slanted
\margin{\uline{NB} $F$ et $G$ \unsure{ens\textsuperscript{t}|ens.} comp. sss \deleted{toute} \deleted{strate}\inserted{multistrate} de $F$ est ens. comp. à \deleted{toute} multistrate de $G$}

b) $\forall$ strate ouverte $U$\inserted{$=X^\circ$} de $F$ et strate ouverte \deleted{$\unsure{U|V}$}$V$\inserted{$={\unsure{Y|Z}}^\circ$} de $G$, on a $U\cap V=\varnothing$ ou $U=V$

b$'$) $\ldotp\ldotp\ldotp\ldotp\ldotp$ \deleted{\unsure{On}} a $\deleted{\unsure{U}}X^\circ\cap\deleted{\unsure{V}}Y^\circ=\varnothing$ ou $X=Y$

c) $\forall\,\unsure{x|X}\in F,Y\in G$, $X\cap Y$ est saturé dans $|F|$ et \unsure{de|dans} $|G|$. \inserted{(\unsure{réentraine|réentraîne} la saturation pour les str. de multistr. \deleted{\unsure{qui}} $\widehat{x},\widehat{y}$)}

d) $F\cup G$ est admissible\unsure{,)|,} et $F$ et $G$ sont des sous-ens\unsure{.} \deleted{\unsure{fer}} \uline{fermés} pour la relation $\subset$

On dit alors que \deleted{\unsure{$G$,}} $F$ et $G$ sont ``ensemblistement compatibles''. Il en résulte que

% slanted
\margin{\uline{NB} Soit $H$ une autre famille admissible d'ens. Pour que $H$ soit ens. comp. à $F$ et $G$ (supposés ens. comp. entre eux), il \unsure{f|s} et s qu'il le soit à $F\cup G$.}
\[
\begin{aligned}
&\inserted{F\cap G\ \text{fermé dans } F \text{ et dans } G}\\
&F,G\ \text{fermés dans } F\cup G
\end{aligned}
\]

et $F\cup G$, en tant qu'ens \inserted{(?)} ordonné est l'\deleted{\unsure{ens.}} \deleted{\unsure{som}} somme \unsure{amalgamée} $\unsure{>|.}\quad F\amalg_{F\cap G}G$,

% slanted
\margin{\uline{NB} Si $F$ et $G$ sont compatibles alors deux sous figures $F'\leq F$ et $G'\leq G$ \unsure{le sont ssi|le sont ss.|le sont aussi} Donc si on a une configuration, l'ens des sous-figures de toutes les figures de la conf. est encore une configuration. On peut aussi \unsure{agrandir|adj.} en y incluant toutes les réunions finies \unsure{d'ens.|de} \deleted{\unsure{telles}} figures \unsure{de la conf.(ig)|de la config} $+$ \unsure{arbitraires|arbitraire}.}

On dit que $F$ et $G$ \uline{sont compatibles}, \deleted{si $F\cup G$ est} s'ils le sont ensemblis\deleted{\unsure{st|t}}tement, et si $F\cup G$ est une figure. \inserted{(Alors $F,G$ sont des sous-figures de $F\cup G$)}

\uline{Ex} Si $F$ et $G$ sont des sous-figures d'une figure $H$, alors ils sont compatibles (car \inserted{ils le sont \unsure{ensembl\textsuperscript{t}.|ensembl.}, et} $F\cup G$ est une sous-figure de $H$, donc une figure.)

Configuration : ensemble de figures deux à deux compatibles.
\scanpage{30}\mspage{29}
\uline{Cont 4} Si $F,\unsure{G|C},$ \deleted{$H$ sont des figures}\inserted{sont \unsure{des|deux}} figures, telles que \unsure{tout|toute} multistrate de $F$ est comp. à \unsure{tout|toute} multistrate de $G^{*}$, alors $F$ et $G$ sont compatibles (l'inverse étant évident)

% slanted
\margin{\uline{NB} \unsure{Alors} $F$ et $G$ sont ens. comp.}

\deleted{\uline{\unsure{Cor}} Si $F,G,H$ sont \unsure{. tout|. trois} $(\unsure{F_i|F_j}\unsure{]|)}$}

\deleted{\unsure{Équi|Éq}} Ceci implique

\uline{Cor} Soit $(F_{\unsure{i|j}})_{i\in I}$ une famille \deleted{\unsure{finie}} de figures. Pour que \deleted{les $F_i$ soient} cette famille soit une configuration, il f et s que pour tout $i\ne j$ et \unsure{tout|toute} multistrate $\widehat{X}\deleted{\unsure{\in}}$ de $F_i$, $\widehat{Y}$ de $F_j$, $\widehat{X}$ et $\widehat{Y}$ soient compatibles. Si $I$ est fini, cela implique que $F=\bigcup F_i$ est une figure, et \inserted{les} $F_i$ \deleted{est} sont des sous-figures.

\deleted{\uline{NB} Inversement, la validité de cet énoncé pour \deleted{\unsure{des}} \unsure{tout|toute} paire $\{F,G\}$ de figures (savoir, \deleted{\unsure{si les}} le critère de}

\uline{Contrées \deleted{ensemblistes}\inserted{de points}} La contrée est dite ``ensembliste'', ou ``\uline{contrée de points}'' (\unsure{ici|là|i.e.} ses lieux prennent le nom de \uline{points}) si deux figures ensemblistement compatibles sont compatibles (ce qui rend donc tautologique Cont 4) 
\scanpage{31}\mspage{30}

Pour les contrées non ensemblistes, il y a lieu de supposer ceci (à vérifier sur les exemples non ensemblistes à venir)

\margin{(?)}
% slanted
\margin{\unsure{remplace} \uline{Cont} 4}
\uline{Cont} $\substack{1\\4}$ Pour que deux figures \inserted{$F,\unsure{G|C}$} soient \margin{\rule{0.35pt}{4.5\baselineskip}}compatibles, il f. et s. qu'elles le soient ensemblistement, et que pour deux \deleted{strates \unsure{dist.}} strates ouvertes $\unsure{U}$ de $F$, $V$ de $G$, $U\ne V$, \deleted{on ait} et pour \deleted{tout} $x\in U$, $y\in V$, $\deleted{\unsure{\{x\}\{\}\{|\{x\}\}}-\{x,y\}}$ $\{\{x\},\{y\}\}$ soit une figure (i.e $\{\{x\}\}$ et $\{\{y\}\}$ soient compatibles

% slanted
\margin{\uline{NB} Les contrées ``\unsure{ensemblistes}'' ou ``contrées de pts'' sont celles qui satisfont $\mathrm{Cont}'_{4}$, et telles que toute partie de $L$ à 2 él. est une 0-sphère i.e. deux lieux \unsure{+} distincts y sont disjoints.}

\uline{NB} Pour tout $x\in L$, on note $\unsure{F|E}_x$ la figure $\{\{x\}\}$, \unsure{consistant} en la seule strate $\{x\}$.

% slanted
\margin{\uline{NB} Un lieu $x$ de $L$ est dit \uline{ponctuel} si si $\forall y\deleted{\unsure{\ne}}\unsure{\in}L$, $y\ne x$, \unsure{$y$} est disjoint de $x$. Ainsi, les contrées de points sont celles où tout lieu est un pt. !}

\deleted{Dire que les fig} Si $x,y\in L$ dire que $F_x$, $F_y$ sont compatibles, signifie que $\{\deleted{\unsure{\{x\},\{y\}}}F_x,F_y\}$ est \deleted{une} une figure. C'est toujours vrai si $x=y$. \deleted{Si $x=y$, \unsure{on dit alors}} \inserted{On dit} que \deleted{\unsure{$\{$}} $x,y$ sont ``\uline{disjoints}'' s'ils sont distincts et compatibles \deleted{\unsure{i.e ssi} \unsure{$\{x$}} On dit aussi \inserted{alors} que $\{x,y\}\in P_2(L)$ est une 0-\uline{sphère} de $C$.

Deux parties \deleted{\unsure{de}} $A,B$ de $L$ sont dit \uline{disjointes} (au sens de la contrée $C$ -- ou $C$-disjointes) si $\forall x\in A$, $\forall y\in B$, $x,y$ sont disjoints ($\Longrightarrow A\cap B=\varnothing$).
\scanpage{32}\mspage{31}

Ainsi $\mathrm{Cont}'_4$ dit que deux figures $F,\unsure{G|C}$ sont compatibles sss pour \unsure{tt|tte} strate ouverte $\deleted{\unsure{A|X}}\unsure{U}$ de $F$ et $V$ de $G$, \deleted{telles que $U\ne V\unsure{_{\!a}|\,a}$, on ait} on ait $U=V$ ou $U$ et $V$ sont $C$-disjoints.

\uline{12. \unsure{Juin|juin}} Depuis hier j'ai pas mal réfléchi à la combinatoire des pseudodisques topologiques. Il est apparu que si $\deleted{\unsure{L|X|D}}$ $(D,\partial D)$ est un tel pseudo-disque (\unsure{en|au} un sens convenable), alors les ``lieux'' sur $D-\partial D$ ne sont \uline{pas} seulement des parties (compactes contractibles etc) de $D-\partial D$, satisfaisant à certaines conditions, mais elles doivent être accompagnées\unsure{.} d'une donnée supplémentaire, une \unsure{classe|donnée} d'\unsure{équivalence|équidistance} de ``\unsure{cordes|corde}'' les joignant à $\partial D$. Du coup, cela montre que l'interprétation \deleted{\unsure{que}} que j'avais en vue, d'une structure de contrée en termes \unsure{dé|de} donnée de l'ens $L$ des lieux, et de certains ensembles de parties de $L$, les ``figures'', n'est pas
\scanpage{33}\mspage{32}

\unsure{:} valable pour la ``topologie combinatoire sans points''. Donc il \unsure{faux|faut} revenir à la donnée de
\diagram{\begin{tikzcd}[ampersand replacement=\&,column sep=0.3em,row sep=1.1em]
C=\bigl(M, \& {\ll,}\arrow[d,out=-90,in=160] \& {R_M,}\arrow[d,out=-90,in=160] \& \text{Fam. de parties \deleted{\unsure{fermées}} de }M\bigr) \\
\& \substack{\text{raffinement}\\\text{(rel d'ordre}} \& \substack{\text{compatibilité}\\\text{\unsure{rel,|rel.} sym. réfl.}} \&
\end{tikzcd}}

% slanted
\margin{$R_M$ résulte de \unsure{``}$\ll,\unsure{\mathcal{F}|\mathscr{F}}$, car \unsure{$\in$|si} \deleted{\unsure{$F,G$|si}} si $X,Y\in M$, $X\,R_M\,Y$ sss $\deleted{\unsure{(X,Y)}}$ $\exists F\in\mathcal F$ telle que $X,Y\in F$ D'autre part, on verra plus bas que si $C$ est finitaire i.e. les figures sont finies \inserted{(\unsure{cas|les} les plus importants !)}, alors $\ll,R_M$ déterminent $\mathcal F$.}

Je voudrais, pour commencer, au lieu de dégager un bon syst. d'axiomes, dégager une terminologie géométrique, en notant (s'il y a lieu) au passage les axiomes qui se seront imposés à l'attention en cours de route.

% slanted
\margin{\uline{NB} Une sous-contrée d'une sous-contrée est une \unsure{of|d} sous-contrée.}

\uline{Sous-contrée} : est définie par un sous-ensemble $M'$ \deleted{\uline{stable} \unsure{p|pr} $\ll$} \inserted{qui est} fermé dans $M$ pour $\ll$, i.e. si $X$ est dans $M'$ tout raffinement de $X$ y est aussi. \deleted{\unsure{Il|Je}} \deleted{\unsure{Elle a} Mais on fera attention que} \inserted{On le munira des relations $\ll_{M'}$ et $R_{M'}$} induites par $\ll_M$ et $R_M$. Mais on fera attention qu'il n'est pas \deleted{évident} \inserted{dit} que les axiomes qu'on va imposer à une \unsure{contrée} soient \unsure{tous} stables par passage à une
